\abstract{
Large language models are typically post-trained using supervised fine-tuning (SFT) and reinforcement learning (RL), yet effectively unifying efficient knowledge injection with robust generalization remains challenging. In this work, we provide a training-dynamics analysis showing that SFT can be interpreted as a special case of policy gradient optimization with an extremely sparse implicit reward and unstable inverse-probability weighting, which together lead to single-path dependency, entropy collapse, and gradient explosion. Motivated by this diagnosis, we propose Group Fine-Tuning (GFT), a unified post-training framework that addresses these intrinsic limitations through two mechanisms: Group Advantage Learning, which constructs diverse response groups and derives normalized contrastive supervision to alleviate reward sparsity, and Dynamic Coefficient Rectification, which adaptively bounds inverse-probability weights to stabilize optimization while preserving efficient knowledge injection. Experiments demonstrate that GFT consistently surpasses SFT-based methods and yields policies that integrate more smoothly with subsequent RL training.
}
